\documentclass[a4paper, 12pt, oneside, BCOR=1cm,toc=chapterentrywithdots]{scrbook}

\usepackage{graphicx}                   % use for pdfLatex
\usepackage{makeidx}                    % \printindex
\usepackage[colorlinks=false]{hyperref}
% \usepackage{tocbibind}
\usepackage{blindtext}
\usepackage{subfigure}
\usepackage{acronym}
\usepackage[german]{babel}
\usepackage{csquotes}
\usepackage[backend=biber,style=alphabetic,backref]{biblatex}
\addbibresource{bibliography.bib}

\defbibfilter{articles}{
  type=article or
  type=inproceedings
}

\hypersetup{bookmarksnumbered=true,bookmarksopen=true, bookmarksopenlevel=1,pdfborder=0 0 0}

% remove this later
\setcounter{tocdepth}{4}

% \defcitealias{ISO924111}{\citefield{ISO924111}{shorttitle}}

\begin{document}
\begin{titlepage}

  \begin{center}
    \raisebox{-1ex}{\includegraphics[scale=1.5]{assets/tuc.pdf}}\\
  \end{center}
  \vspace{0.5cm}

  \begin{center}

    \LARGE{\textbf{Title of the bachelor thesis}}\\
    \vspace{1cm}


    \Large{\textbf{Bachelor Thesis}}\\
    \vspace{1cm}
    Submitted in Fulfilment of the\\
    Requirements for the Academic Degree\\
    B.Sc.\\
    \vspace{0.5cm}
    Dept. of Computer Science\\
    Chair of Computer Engineering
  \end{center}
  \vspace{3cm}
  Eingereicht von: Joshua Jeschek\\
  Matrikelnummer: 652272\\
  Datum: 24.12.2023\\
  \vspace{0.3cm}\\
  Dr. Thomas Wilhelm-Stein \\
  Dr. Heino Rudolf

\end{titlepage}

\addchap*{Abstract}
\blindtext


\tableofcontents

\begingroup
\cleardoublepage
\addcontentsline{toc}{chapter}{\listfigurename}
\listoffigures

\let\clearpage\relax

\cleardoublepage
\addcontentsline{toc}{chapter}{\listtablename}
\listoftables
\endgroup

\twocolumn
% \addchap{List of Abbreviations}
% \begin{acronym}[Bash]
% 	\acro{SQL}{Structured Query Language}
% \end{acronym}

\onecolumn

\chapter{Einleitung}

\section{Hintergrund und Motivation}
\section{Ziele und Umfang}
\section{Aufbau der Arbeit}

\chapter{Theoretische Grundlagen}

\section{envVisio}
\subsection{Das Realitätsmodell}
\subsubsection{Allgemeine Erklärung von envVisio}
\subsubsection{Vielleicht Aufteilung in mehrere Abschnitte?}
\subsubsection{was ist alles wichtig?}
\subsection{Laden und Konvertieren}
\subsubsection{auf Loader und insbesondere Converter eingehen}
\subsubsection{da in Converter Formular eingesetzt werden soll}
\subsubsection{wichtig, Prozesse zu verstehen}

\section{Human-Centered Design}
\subsubsection{Begriffe und Konzepte aus Design of Everyday Things}
\subsubsection{Abschnittstitel nicht final, vielleicht is HCD auch ein Unterpunkt auf gleicher Höhe wie Aff., Sign. und Map.}
\subsection{Affordances, Signifiers und Mappings}
\subsubsection{Kann gut bei 2.3.2 angewendet werden}
\subsection{Activity-Centered Design}

\section{Formulare}
\cite{wroblewskiWebForm2008}
\subsection{Funktionsweise}
\subsubsection{Technische Fkt.-weise von F.}
\subsection{Designanfordungen}
\cite{normanDesignEveryday2013} (auf Affordances, Signifiers und Mappings eingehen)
\subsubsection{Rückbezug auf 2.2}
\subsubsection{Anordnung}
\subsubsection{Kontrollelemente}
\subsubsection{Bezeichnungen und Platzhalter}
\subsubsection{Validierung}

\clearpage
\section{Usability Testing}
\section{Usability Testing}

Um auf \textit{Usability Testing} einzugehen, lohnt es sich, einen genaueren Blick auf \textit{Usability} zu werfen. \citefield{ISO924111}{shorttitle} definiert \textit{Usability} wie folgt:
\begin{quote}
  extent to which a system, product or service can be used by specified users to achieve specified goals with effectiveness, efficiency and satisfaction in a specified context of use
  \hspace*{\fill{}}\citetalias{ISO924111}
\end{quote}
Diese Definition ist kurz, umfasst jedoch trotzdem drei Hauptelemente, die \textit{Usability} ausmachen. Es sollte möglich sein, Aufgaben mit ausreichender Geschwindigkeit auszuführen (\textit{efficiency}) und das erwünschte Ziel zu erreichen (\textit{effectiveness}). Der dritte wichtige Faktor ist die Anwenderzufriedenheit (\textit{satisfaction}), welche mehr auf die subjektive Erfahrung von Nutzer:innen eingeht. So kann eine Aufgabe schnell und korrekt ausgeführt werden, aber dennoch nicht zu Zufriedenheit führen, da beispielsweise wichtige Informationen nicht angezeigt wurden, die über den Erfolg informieren. Der gleiche Effekt kann auch andersherum stattfinden - sollte die persönliche Einschätzung von Zufriedenheit sehr gut sein, können womögliche Probleme der Effizienz und Effektivität vernachlässigt werden, sodass das Produkt trotzdem als gut eingeschätzt wird.
\parencite{barnumUsabilityTesting2021}

Des Weiteren betont \citeauthor{barnumUsabilityTesting2021}, dass es sich um spezifische Nutzer:innen, Ziele und Nutzungskontexte handelt. Das bedeutet, dass \textit{Usability} nur für eine spezifische Gruppe an Menschen betrachtet werden kann, welche das Produkt benutzen werden und deren Ziele denen des Produkts entsprechen müssen. Außerdem kann die \textit{Usability} auch nur für den angedachten Nutzungskontext - die Umgebung in der ein Produkt genutzt wird - betrachtet werden.
\parencite{barnumUsabilityTesting2021}

\textcite{quesenberyDimensionsUsability2003} und \textcite{nielsenUsabilityEngineering1994} spezifizieren folgende 5 Attribute der \textit{Usability}:
\begin{itemize}
  \item \textbf{Effektivität}: wie umfassend und akkurat Aufgaben abgeschlossen werden,
  \item \textbf{Effizienz}: wie schnell dies erfolgt,
  \item \textbf{Ansprechende Gestaltung}: wie gut Interaktionen gesteuert werden und für Zufriedenheit gesorgt wird (von \citeauthor{nielsenUsabilityEngineering1994} vernachlässigt)
  \item \textbf{Fehlertoleranz}: wie das System Fehler verhindert und sich von Fehler erholt
  \item \textbf{Erlernbarkeit}: wie intuitiv das System ist, und wie Lernen im Laufe der Nutzung ermöglicht wird (von \citeauthor{nielsenUsabilityEngineering1994} in zwei Attribute geteilt)
\end{itemize}
Laut \textcite{barnumUsabilityTesting2021} hat der Faktor der Zufriedenheit in den letzten Jahren an Bedeutung zugenommen. Das wird auch davon unterstützt, dass \textcite{nielsenUsabilityEngineering1994} weniger Wert auf eine ansprechende Gestaltung legte als \textcite{quesenberyDimensionsUsability2003}.

\newpage

\textit{Usability Testing} ist ein Sammelbegriff für Methoden bei denen Benutzer:innen mit einem System interagieren, um ein Ziel zu erreichen oder Szenario umzusetzen. Dabei sind die Umstände der Interaktionen kontrolliert und Verhaltensdaten werden gesammelt \parencite{wichanskyUsabilityTesting2000}. Ziel ist es, die \textit{Usability} des Systems zu überprüfen und eventuelle Probleme zu identifizieren.

Dabei ist es laut \textcite{barnumUsabilityTesting2021} wichtig, dass die Nutzer:innen Aufgaben erfüllen, die als echt empfunden und von Bedeutung für sie sind. So werden Probleme, die bei der Nutzung in der realen Welt auftreten effizienter aufgedeckt.

Wichtige Aspekte des \textit{Usability Testing} werden im weiteren Verlaufe des dieses Kapitels beleuchtet.

\subsection{Testmethoden}
\subsubsection{Quantitative vs Qualitative Tests}
\subsubsection{Formative vs Summative Studien}
\subsection{Teilnehmer:innen}
\subsubsection{Wer wird ausgewählt?}
\subsection{Aufgaben und Scenarios}
\subsection{Evaluation}
\subsubsection{Wie können Fragebögen zur Evaluierung erstellt werden?}
\subsubsection{Verwendung von Standard-Fragebogen}

\section{Analyse und Vergleich von existierenden Projekten}
\subsection{Home Assistant}
\subsection{DBSnap}
\url{https://www.semanticscholar.org/paper/DBSnap\%3A-Learning-Database-Queries-by-Snapping-Silva-Chon/aabf7791575f8cc62432ca5e6289862eec0f6845}
\url{https://www.semanticscholar.org/paper/DBSnap-2\%3A-New-Features-to-Construct-Database-by-Silva-Loza/61c95757d7d4151f51ea848623c6839816756385}
\subsection{SQheLper}
\url{https://www.semanticscholar.org/paper/SQheLper\%3A-A-block-based-syntax-support-for-SQL-Jacobs-Jaschke/b4f6023a4ee72d0103ee521c4f9b35768d38810b}
\subsection{Scratch}


\chapter{Implementierung}
\section{Zielformulierung}
\subsubsection{Genaue Formulierung der Problemstellung, was ist das Ziel?}
\section{Lösungsvorstellung}
\subsubsection{Erklärung des Resultats}

\chapter{Studie}
\section{Studienaufbau}
\subsection{Problemformulierung und Ziele}
\subsection{Auswahl von Teilnehmer:innen}
\subsection{Aufgaben}
\subsection{Scenarios}
\subsection{Evaluation}
\subsection{Studienablauf}
\section{Ergebnisse und Analyse}
\section{Schlussfolgerungen und Zukunftsaussicht}

\chapter{Fazit}


\nocite{*}
\chapter*{Quellen}
\begingroup
\let\clearpage\relax
\printbibliography[filter=articles, title=Artikel, heading=subbibliography]
\printbibliography[type=book, title=Bücher, heading=subbibliography]
\printbibliography[type=misc, title=sonstige, heading=subbibliography]
\endgroup

\printindex

\end{document}
